\documentclass{article}

\usepackage{times}
\usepackage{hyperref}
\usepackage[cm]{fullpage}
\usepackage{fancyhdr}
\usepackage{graphicx}
\usepackage{natbib}

%\setlength{\headheight}{9.2pt}
\setlength{\headsep}{5.2pt}
\pagestyle{fancyplain}
\pagenumbering{gobble}
\lhead{\textbf{{\large Toby Dylan Hocking}  \\ \texttt{toby.hocking@mail.mcgill.ca} }}

\rhead{\includegraphics[height=0.75cm]{logo-mcgill}}

\begin{document}

 \mbox{ }
% \vspace{ -0.5cm}

\section*{\centering Research Statement}

During 2009-2012, I completed my PhD degree in applied mathematics
from the Ecole Normale Sup\'erieure de Cachan, France. I studied
machine learning algorithms for biology and medicine, under the
supervision of Dr. Francis Bach (INRIA, ENS) and Dr. Jean-Philippe
Vert (Institut Curie, Mines ParisTech). I am currently continuing my
research about machine learning algorithms for biology, as a
postdoctoral research associate under the supervision of Dr. Guillaume
Bourque, in the McGill University Department of Human Genetics.

With recent advances in technology, more data are being generated and
recorded than ever before in human history. For example, in scientific
fields such as medicine and biology, new sensors and DNA sequencing
machines are being used to generate many high-dimensional data
sets. This has led to a need to develop new statistical models and
machine learning algorithms to understand patterns in these data, i.e.
to generate knowledge about how the human genome influences medical
conditions. I am interested in developing new statistical machine
learning algorithms to address these needs, by collaborating with
domain experts who generate these large data sets.

So far my research has been focused on proposing new supervised and
unsupervised machine learning models for clustering, regression,
changepoint detection, ranking, and classification. My contributions
to the statistics and machine learning literature are algorithms for
efficiently solving these discrete and convex optimization problems,
which typically involve very large data sets. I have developed close
collaborations with domain experts in the fields of medicine and
biology, in order to better understand their problems and propose
appropriate solutions. The following sections summarize my previous
research projects, and my future research plan.

\section{Changepoint detection and regression algorithms for DNA copy
  number data}

During my PhD in Paris, I worked in collaboration with medical doctors
at the Institut Curie on semi-automatic methods for diagnosing
neuroblastoma, a cancer that most frequently occurs in small
children. The doctors wanted to give targeted treatments based on the
number of abrupt changes in DNA copy number data
from the patient's tumor. Previous unsupervised algorithms for
changepoint detection in these data provided insufficient accuracy
with respect to the doctors' interpretation of the signal/noise in
these data, and my goal was to develop new algorithms for more
accurate changepoint detection.

\paragraph{New labeling method for supervised changepoint detection.}

In work published in \emph{BMC Bioinformatics}
\citep{HOCKING-breakpoints}, we proposed to re-frame the changepoint
detection problem in terms of supervised learning. In our framework,
the doctors first labeled a subset of their data with positive regions
that contain changepoints, and negative regions that contain no
changepoints. We thus have a set of labeled data sequences, and the
goal is to learn a model that provides accurate changepoint predictions
on a set of un-labeled data sequences. The main contributions of our
paper were to show that existing algorithms can be (1) trained by
learning parameters that minimize the number of
incorrect training labels, and (2) compared by computing the number of
incorrect test labels in a held-out data set. We performed a
comparison study which showed that the most accurate existing method
was a Gaussian maximum likelihood changepoint detection model. 
%This result was somewhat expected, because this model is the sparsest
% Each
% data sequence is modeled by a piecewise constant segment mean, and a
% dynamic programming algorithm finds the most likely changepoint
% positions, subject to an L0 sparsity constraint on the number of
% changes. 
In this paper we selected the number of changes using a
simple univariate penalty function that depends on the number of data
points, and in our next paper we investigated generalizing that to
multivariate penalty functions.

\paragraph{Regression models for supervised penalty function
  learning.}

In work published at \emph{ICML'13} \citep{HOCKING-penalties}, we
proposed a fast and accurate supervised machine learning algorithm
that exploits the structure of the labeled changepoint detection
problem. We showed that learning a penalty function for selecting the
number of changepoints is in fact a regression problem with censored
outputs. For each data sequence in which we would like to detect
changepoints, the input is a feature vector of statistics (number of
data points in the sequence, variance estimates, quantiles, etc), and
the output is an interval of log(penalty) values that achieves the
minimum number of incorrect labels. We showed that a data sequence
with a negative label results in a right-censored output (when the
penalty is too small, there are too many changepoints); a positive
label results in a left-censored output (when the penalty is too
large, there are not enough changepoints). When a data sequence has
both positive and negative labels, the output is interval-censored.
We proposed a margin-based convex loss function that exploits the
structure of the censored outputs, and an L1 penalty to regularize and
encourage a sparse penalty function. Our main contribution was a new
convex optimization algorithm that minimizes the resulting convex
objective function. We showed that our algorithm learns a multivariate
linear penalty function, and provides more accurate changepoint
detection than the previous univariate penalty. I implemented these
algorithms in the penaltyLearning R package, which I presented during
a useR2017 tutorial on optimal changepoint detection algorithms
\citep{change-tutorial}. In more recent work under review at
\emph{NIPS'17}, we proposed a nonlinear decision tree model for
learning such penalty functions. 

\paragraph{Collaborations using an interactive web app for supervised
  changepoint detection.} In work published in \emph{Bioinformatics}
\citep{hocking-SegAnnDB}, we proposed an web application that
facilitates interactive labeling and supervised machine learning
analysis of DNA copy number data. On the practical side, our system
was novel because it allows genomic scientists to not only view the
data but also label obvious signal and noise patterns. This
functionality is essential for collaborations, because when a domain
expert sees an incorrect prediction, he/she will be able to
immediately correct it by adding appropriate labels. Our main
statistical contribution was an efficient dynamic programming
algorithm for computing the most likely K changepoints for a given set
of K positive labels. This algorithm was essential in order to provide a model
which fits any set of labels provided by an expert user. We have used
this system to facilitate collaborations with scientists at Intitute
Curie, Paris, France \citep{Chicard}; and Aichi Cancer Center, Nagoya,
Japan \citep{Hocking-Leukemia-2016,m14:clonal}.


\paragraph{Pruning algorithms for optimal changepoint detection.} In
work published in \emph{Statistics and Computing} \citep{fpop}, we
studied two previous optimal changepoint detection algorithms, and
generalized their pruning techniques to other optimization
problems. One previous algorithm used an inequality pruning technique
for the penalized optimization problem, and we showed that a
functional pruning technique can also be used to solve the penalized
optimization problem. Our mathematical analysis showed that every
changepoint pruned by the inequality technique will also be pruned by
the functional technique. This theoretical result suggested that the functional
pruning algorithm should be faster, and we observed that it is indeed
faster in an empirical analysis of the algorithm on several real DNA
copy number data sets from neuroblastoma tumors. My main contribution
to this paper was to conduct the empirical analysis of speed and
accuracy of the algorithm. In a comparison with other changepoint
detection algorithms, we showed that our proposed algorithm achieves
state-of-the-art speed and changepoint detection accuracy.

\section{Peak detection in ChIP-seq data}

\paragraph{New labeling method for supervised peak detection.} In work
published in \emph{Bioinformatics} \citep{HOCKING-chipseq}, we
proposed a labeling method and supervised learning framework for peak
detection in genomic data such as ChIP-seq. These experiments are
performed in order to determine active and inactive regions of the
genome in different individuals and cell types. In these data, the
main analysis task is to segment the genome into peaks (with large
values) and background noise (with small values). Although there are
dozens of unsupervised heuristic algorithms for this task in the
bioinformatics literature, they all suffer from inaccurate peak
detection because there is no method to systematically adapt tuning
parameter values to the different peak patterns in different data
sets.

We proposed to frame the problem in terms of supervised learning, with
labels from expert scientists that indicate regions in specific
samples with or without peaks. We showed that these labels can be used
to (1) choose appropriate parameters that minimize the number of
incorrect labels in training data, and (2) compare the peak detection
accuracy of different algorithms in held-out test data. In our
cross-validation experiments, we trained a single significance
threshold parameter of each unsupervised algorithm, then computed
predicted peaks and number of incorrectly predicted labels in a
held-out test set. We observed that different algorithms achieved
maximum peak detection accuracy in different data sets; there was no
single algorithm which was most accurate in all data sets. We also
showed that labels from different scientists are highly consistent, by
observing similar test error rates when training on labels from the
same or a different scientist.

\paragraph{Up-down constrained changepoint detection model for peak
  detection.} In work published at \emph{ICML'15}
\citep{HOCKING-PeakSeg}, we proposed a constrained changepoint
detection model for peak detection in ChIP-seq data. We observed that
the existing optimal changepoint model could sometimes be interpreted
in terms of peaks (after up changes) and background (after down
changes). However, in some data sets the maximum likelihood model
consists of several consecutive up changes, which is not directly
interpretable in terms of peaks and background. We thus proposed an
additional constraint that up changes must be followed by down
changes, which results in a more interpretable model (with lower
likelihood). Our main contribution in this paper was to adapt the
classical quadratic time dynamic programming algorithm to this new
constraint. Another contribution was to demonstrate that we can
efficiently learn a penalty function, using the same censored
regression algorithm that we previously used to predict the number of
changes in the unconstrained changepoint detection model. We showed
that this method achieves state-of-the-art peak detection accuracy in
all of the labeled data sets that we created, thereby proving that a
single algorithm can adapt to peak patterns in different data sets.

Our quadratic time dynamic programming algorithm has two issues: it is
slow for large data sets, and it does not necessarily compute the
optimal solution to the constrained problem. In work under review in
\emph{Annals of Applied Statistics}
\citep{Hocking-constrained-changepoint-detection}, we propose a new
algorithm that resolves both issues. We describe how a functional
pruning technique can be generalized to changepoint problems with
constraints between adjacent segment means (e.g. second segment mean
must be greater than first segment mean). This results in a fast
log-linear time algorithm, because it immediately prunes changements
which result in a sub-optimal cost. The algorithm is also guaranteed
to compute the solution to the constrained changepoint detection
problem. We applied this algorithm to our benchmark of labeled
ChIP-seq data sets. We observed the same state-of-the art peak
detection accuracy, and orders of magnitude faster run-times.

\paragraph{Changepoint model for comparing samples.}
Rather than modeling the peaks in each sample in isolation, scientists
often would like to characterize the differences between samples. In
my ongoing collaborations at the McGill genome center, we are trying
to characterize active and inactive regions in several different human
cell types (B cells, T cell, monocytes, etc). In ongoing work, I have
proposed a model that enforces the constraint that peaks should occur
in the same positions across samples \citep{HOCKING-PeakSegJoint}.

\section{Other research projects}

\paragraph{Convex clustering.} In work published at \emph{ICML'11}
\citep{HOCKING-clusterpath}, I described efficient algorithms for
hierarchical clustering using convex fusion penalties. This paper has
been influential, with over 80 citations including theoretical,
algorithmic, and applications papers. 

\paragraph{Bayesian model of evolution.} In work that we published in
\emph{PLOS ONE} \citep{HOCKING-evolution}, we proposed to identify
genetic markers under selection using a hierarchical Bayesian model of
evolution.

\paragraph{The inlinedocs R package for documentation generation} was
described in our \emph{Journal of Statistical Software}
\citep{hocking13:inlinedocs} article.


\paragraph{Support vector machines for ranking and comparing} In my
2013 postdoc at Tokyo Institute of Technology, I studied a support
vector machine algorithm for ranking and comparing
\citep{svmcompare}. In ongoing work which we plan to publish in
\emph{Journal of Machine Learning Research}, we train this
algorithm on historical data of chess matches, and use it to predict
player rankings and future game outcomes.

\section{Plan for future research}

\paragraph{Aim 1: survival models for penalty function learning.} In
previous work I have showed that learning a penalty function for
changepoint detection can be cast as a regression problem with
censored outputs, a problem that has been extensively studied in the
survival literature. In this project, I will use regularized Accelerated Failure
Time models from the survival literature to learn a multivariate
linear penalty function. The main difference with my previous work is
that these models assume a parametric distribution for the penalty values
(e.g. maximum Gaussian likelihood), whereas my previous
work used a discriminative model (minimizing a margin-based loss
function).

\paragraph{Aim 2: multi-task penalty function learning.} My previous
work has shown that labels from a single genomic scientist can be used to
learn a penalty function for changepoint detection. But in practice I
often have several different genomic scientists that have provided
separate labels for a given data set. One could imagine that some
scientists would provide similar labels, and others have different
labels, for the same data sequence. In this setting, I would like to
learn a different penalty function for each scientist, under the
assumption that there will be clusters of similar scientists. In this
project I will investigate adapting techniques from the multi-task
learning literature to this more complex penalty function learning problem.


\paragraph{Aim 3: functional pruning algorithms for multi-state
  changepoint models.} My previous research has shown that a
functional pruning technique can be used to obtain fast and exact
constrained changepoint detection algorithms. In this project, I will
investigate generalizing this technique to more complex constrained
changepoint models than the two-state peak/background model we used
for ChIP-seq data. For example, neuro spike train data sets could be
modeled using a changepoint model with three states: flat/constant,
increasing toward a spike, decreasing after a spike.

\paragraph{Aim 4: GenomicLearner platform for collaborations.} My previous
research has shown that genomic scientists can label some obvious
signal/noise patterns in their data sets, and these labels can be used
to train supervised learning algorithms for highly accurate pattern
recognition. However, existing genomic data visualization software
(e.g. UCSC genome browser) does not support interactive labeling. The
goal of this project will be to develop GenomicLearner, a extensible
web app for machine learning analysis of genomic data. The app will be
used to (1) view and label genomic data sets (e.g. ChIP-seq, copy
number, etc), and (2) interactively run state-of-the-art machine
learning algorithms. GenomicLearner will make recent developments in
machine learning more accessible to genomic scientists, thus
facilitating collaborations between researchers in machine learning
and biology. I have already written a grant proposal, which I plan to
submit to NIH or NSF.

% \paragraph{Aim 3: other labeling and supervised learning in data
%   sequences.}
% My collaborations with genomic scientists have revealed that
% algorithms. Multivariate patterns in histone marks (chromatin
% states). Labeling confidence bands.

\paragraph{Aim 5: Deep learning for nonparametric changepoint
  detection.} My previous work on supervised changepoint detection has
been limited to learning the penalty function that predicts the number
of changes in maximum likelihood segmentation models. These models
work quite well in relatively small data sets (thousands of labels),
but I expect they will not be flexible enough to learn accurate models
in the larger data sets that I will generate using the GenomicLearner
system (millions of labels). In this project I will explore training
deep convolutional neural networks, which should provide more accurate
predictions in these larger data sets.

\bibliographystyle{plain}
\bibliography{TDH-refs}

\end{document}
